\chapter{Subconjunctival Hemorrhage}

\section*{Definition}
Subconjunctival hemorrhage (SCH) is the extravasation of blood into the potential space between the conjunctiva and the underlying episclera, appearing as a sharply demarcated, homogeneous red patch on the ocular surface.

\section*{Pathophysiology}
Bleeding occurs from rupture of fragile conjunctival or episcleral blood vessels due to sudden increased venous pressure, direct trauma, or vessel fragility. The blood collects beneath the transparent conjunctiva and is slowly resorbed over 7--14 days without sequalae.

\section*{Etiology}
\begin{itemize}
  \item \textbf{Valsalva maneuvers:} Heavy lifting, violent coughing, forceful sneezing, or straining with defecation
  \item \textbf{Trauma:} Direct eye rub, foreign body, or blunt ocular injury
  \item \textbf{Systemic hypertension:} Acute blood pressure spikes
  \item \textbf{Anticoagulant therapy:} Aspirin, warfarin, DOACs (apixaban, rivaroxaban), or clopidogrel
  \item \textbf{Coagulopathy:} Thrombocytopenia, hemophilia, von Willebrand disease
  \item \textbf{Ocular infection/inflammation:} Severe conjunctivitis, adenoviral keratoconjunctivitis
  \item \textbf{Idiopathic:} Spontaneous, often upon waking
\end{itemize}

\section*{Indications for Evaluation}
Seek care if:
\begin{itemize}
  \item SCH associated with eye pain, photophobia, or decreased vision
  \item Recurrent or bilateral SCH without clear precipitating cause
  \item Patient on anticoagulation with concern for systemic bleeding risk
  \item SCH accompanied by epistaxis, easy bruising, or gingival bleeding (suspect coagulopathy)
  \item Failure to resolve within 2--3 weeks
\end{itemize}

\section*{Related Symptoms}
\begin{itemize}
  \item Visible red patch on the sclera with clear edges
  \item Mild irritation or foreign body sensation (usually absent)
  \item No pain, discharge, or photophobia (if uncomplicated)
  \item Eyelid ecchymosis if associated with trauma
  \item Hemorrhagic chemosis if extensive
\end{itemize}
\textbf{Note:} A painless, asymptomatic SCH that does not extend beyond the limbus is almost always benign; however, SCH that tracks circumferentially around the cornea suggests a retrobulbar hemorrhage and requires urgent evaluation.

\section*{Self-Care / Home Management}
\begin{itemize}
  \item Cold compresses for the first 24 hours to reduce further bleeding
  \item Warm compresses after 48 hours to accelerate resorption
  \item Artificial tears for mild irritation if needed
  \item Avoid aspirin or NSAIDs if not medically necessary
  \item Reassurance that the blood will resorb in 1--2 weeks, changing color from red to yellow before clearing
  \item Avoid heavy lifting and straining until resolved
\end{itemize}

\section*{Diagnostic Approach}
\begin{itemize}
  \item Slit-lamp examination to confirm blood is beneath the conjunctiva (not the cornea) and posterior extent
  \item Measurement of blood pressure to rule out systemic hypertension
  \item Assessment for evidence of head trauma or orbital fracture if indicated
  \item Coagulation studies (PT, aPTT, INR, platelet count) if recurrent, bilateral, or patient on anticoagulation
  \item Dilated fundus examination to exclude intraocular extension (rare)
\end{itemize}

\section*{Treatment / Management Overview}
\begin{itemize}
  \item No specific treatment required for uncomplicated SCH
  \item Avoidance of precipitating activities until resolution
  \item Management of blood pressure if hypertensive
  \item Adjustment of anticoagulation only in consultation with the prescribing physician
  \item Hematology referral for suspected coagulopathy
  \item Surgical evacuation rarely indicated; only for massive SCH causing extreme proptosis or globe compression
\end{itemize}

\clearpage
